\chapter{绪论}\label{Chap:chap1}

\section{研究背景与动机}

随着云计算、大数据、人工智能、移动互联网和物联网技术的蓬勃发展，端到端网络传输的服务对象正在从固定主机扩展到手机、笔记本、台式机、边缘节点和传感设备等多模态终端，传输路径也从近端局域链路延伸到跨城、跨域、蜂窝接入、无线局域网和卫星通信等远距离复杂网络。传输控制协议（Transmission Control Protocol，TCP）作为互联网传输层的核心协议，承载了大量关键业务流量，其拥塞控制机制的性能直接影响网络资源利用效率、应用响应延迟和用户服务质量。网络拥塞控制是计算机网络领域最具挑战性的研究课题之一，自互联网诞生以来一直是学术界和工业界持续关注的焦点问题。

在现代端到端网络中，链路与终端条件呈现出更强的异构性。一方面，城域/广域网、蜂窝网络、WiFi与卫星链路具有不同的RTT量级、丢包背景、可用带宽和链路抖动；另一方面，手机、电脑、边缘设备等终端在处理能力、无线接入方式、移动性和业务负载上差异显著。这种长距离、多接入、多终端并存的网络环境对传统拥塞控制算法提出了前所未有的严峻挑战。

现代网络应用承载着多种类型的业务负载，这些业务对网络性能有着截然不同的需求。移动视频、实时会议和云游戏需要稳定的低延迟与低抖动；分布式存储、内容分发和大文件同步需要高吞吐量以支持大规模数据传输；边缘计算、工业物联网和车联网则同时关注可靠性、实时性和终端资源约束。单一的拥塞控制策略难以同时满足这些多样化的性能需求。

更为复杂的是，远距离和多接入传输中的拥塞信号更难解释。丢包可能来自真实队列溢出，也可能来自无线误码、链路切换或终端侧调度抖动；RTT升高既可能表示排队，也可能来自路径变化、接入链路波动或终端处理延迟；移动终端后台同步、视频传输、RPC突发、备份任务和多租户共享链路还会造成短时间带宽挤占。传统单一信号驱动的拥塞控制方法难以在这些场景中稳定区分拥塞程度。

传统的TCP拥塞控制算法在这种复杂的网络环境中显露出诸多不足。基于丢包的算法（如TCP NewReno\upcite{newreno}、TCP CUBIC\upcite{cubic}）将数据包丢失作为拥塞信号，但丢包是拥塞的滞后指示器，通常意味着网络缓冲区已经溢出，排队延迟已经积累到较高水平。基于延迟的算法（如TCP Vegas\upcite{vegas}、TCP BBR\upcite{bbr}）尝试通过监测RTT变化提前感知拥塞，但RTT测量易受多种因素干扰，且在与基于丢包算法共存时存在公平性问题。

显式拥塞通知（Explicit Congestion Notification，ECN）机制为解决上述问题提供了新的思路。ECN允许网络设备在队列长度超过阈值时标记数据包而非丢弃，从而实现更早、更精确的拥塞信号传递。数据中心TCP（DCTCP\upcite{dctcp}）验证了ECN细粒度反馈在特定低时延网络中的有效性，但其适用前提较强；在远距离、多接入和多模态终端并存的环境中，现有算法对ECN、RTT、丢包和协议栈状态的联合利用仍不充分，大多采用静态参数配置，缺乏对动态网络环境的自适应能力。

近年来，机器学习特别是强化学习技术为网络拥塞控制提供了新的研究思路\upcite{rlcc}：Aurora、Orca等方案通过训练神经网络策略或在线学习调整发送行为，理论上能够自动适应网络环境。然而，现有学习型方案仍存在诸多工程障碍：状态空间设计不完善，通常只使用4-6维的有限状态参数，未利用ECN与拥塞控制状态机等协议栈内部信号；需要离线训练与模型推理设施，训练分布之外的泛化能力有限；决策过程不透明，难以逐事件审计与复现。这些问题限制了学习型方案在异构网络环境中的实际应用效果。

本文因此选择\emph{白盒启发式}的设计路线：不引入学习型策略，而是在协议栈回调点采集包含ECN、拥塞状态机与窗口/RTT在内的多维状态，以确定性规则在每个确认事件上直接输出拥塞窗口与慢启动阈值决策，仅在参数层使用基线相对的轻量反馈自适应。该路线保留了传统算法低开销、可解释、可复现的优点，同时通过多信号融合与参数自适应获得跨场景的稳定性能。

\section{研究目标与主要贡献}

针对上述问题，本文提出了Swift算法——一种启发式多信号融合自适应网络拥塞控制方法。"Swift"意为"雨燕"，雨燕以其轻盈敏捷的飞行姿态和清脆悦耳的歌声著称。选用"Swift"作为算法名称，寓意本算法追求的设计目标：如同雨燕般轻盈敏捷，实现高吞吐量（High Throughput）与低延迟（Low Latency）的完美平衡；如同雨燕对环境变化的敏锐感知，实现对多种拥塞信号的快速响应和智能融合。

图\ref{fig:swift_contribution_overview}展示了Swift算法的三项核心贡献及各组件之间的逻辑关系。

\begin{figure}[htbp]
\centering
\begin{tikzpicture}[
    node distance=1.0cm and 1.2cm,
    block/.style={rectangle, draw, fill=blue!8, text width=4.6cm, minimum height=1.0cm, align=center, rounded corners=3pt, font=\small},
    arrow/.style={->, >=stealth, thick, draw=gray!70},
]
% 顶层目标
\node[block, fill=red!12, text width=10cm, minimum height=0.9cm] (goal) {Swift设计目标：面向远距离传输与多模态终端的高吞吐量 $+$ 低延迟 $+$ 稳定性};

% 第二层：三项核心创新
\node[block, below left=1.0cm and -0.2cm of goal] (obs) {贡献1：多信号状态感知\\{\scriptsize 11维有效观测子空间，纳入ECN、CA状态、RTT和在途字节}};
\node[block, below=1.0cm of goal] (detect) {贡献2：拥塞判定与差异化响应\\{\scriptsize 丢包/ECN/超时分级，连续递减保护}};
\node[block, below right=1.0cm and -0.2cm of goal] (fusion) {贡献3：启发式反馈自适应融合控制\\{\scriptsize RTT膨胀、反馈值基线比较、BDP估计与窗口有界控制}};

% 底层
\node[block, below=1.2cm of detect, fill=green!8, text width=10cm] (impl) {实现验证：ns-3 $+$ ZeroMQ同步信道 $+$ 多场景对照实验};

% 箭头
\draw[arrow] (goal.south) -- ++(0,-0.25) -| (obs.north);
\draw[arrow] (goal.south) -- (detect.north);
\draw[arrow] (goal.south) -- ++(0,-0.25) -| (fusion.north);
\draw[arrow] (obs.south) -- ++(0,-0.35) -| (impl.north);
\draw[arrow] (detect) -- (impl);
\draw[arrow] (fusion.south) -- ++(0,-0.35) -| (impl.north);
\end{tikzpicture}
\caption{Swift算法核心贡献概览}
\label{fig:swift_contribution_overview}
\end{figure}

本文的主要研究贡献如下：

\textbf{（1）面向远距离与异构终端的多信号状态感知}

本文设计了一个11维有效观测子空间（嵌入于15元素状态传输容器，另含套接字UUID、环境类型、仿真时间、节点ID四项跨进程路由用的元数据通道），将显式拥塞通知（ECN）状态、拥塞控制状态机状态、拥塞事件类型、RTT和在途字节等TCP协议栈信息纳入决策状态表示。与现有方案通常采用的4--6维简化状态空间相比，Swift的观测空间能够为手机、电脑、边缘节点和长距离链路上的TCP连接提供更全面、更精细的状态感知能力，为精准的拥塞控制决策奠定基础。

\textbf{（2）多信号融合拥塞判定与差异化响应}

本文提出了一种多信号融合的拥塞判定方法，采用分层优先级机制综合分析丢包、超时、ECN和拥塞控制状态信号。该方法的核心思想是：不同类型的拥塞信号具有不同的可靠性、时序特征和严重程度，应当采用差异化处理策略。具体而言，丢包和超时代表较严重的拥塞或路径异常，ECN标记代表队列溢出前的早期预警，拥塞控制状态机则提供协议栈内部语义。基于该判定结果，Swift针对丢包、ECN和超时分别采用0.70、0.75和0.50的窗口保留因子，并引入连续缩减保护机制，限制最多3次连续缩减后保持窗口不变，避免因连续拥塞事件导致窗口过度收缩。

\textbf{（3）启发式反馈自适应的融合控制}

本文实现了启发式反馈自适应的融合控制策略，根据RTT膨胀比、性能反馈值相对慢速基线的指数移动平均偏离和连续增长趋势动态调整乘性增加因子$\alpha$，并将其限制在$[0.85, 1.30]$范围内以保证控制稳定性。在窗口计算中，Swift采用时间窗口投递速率采样加窗口化最大值过滤的两级方法估计BDP，并围绕$\alpha \times BDP$目标窗口进行有界步长的渐进拉升或回落；同时结合RTT约束下的管道利用率感知加性递增机制，在吞吐量、延迟和稳定性之间实现自适应折中。上述策略在ns-3仿真平台与ZeroMQ进程间同步信道上实现，并通过多场景对照实验完成验证。

\section{论文组织结构}

本文的其余部分组织如下：

第2章介绍相关工作，包括传统拥塞控制算法、基于机器学习的拥塞控制方法、显式拥塞通知机制、智能决策方法在网络领域的应用综述等。

第3章详细描述Swift算法的设计与实现，包括系统架构、观测空间设计、多信号融合拥塞检测、差异化拥塞响应、自适应参数调优、融合控制策略等核心组件。

第4章介绍实验评估，包括实验环境配置、评测指标、对比算法、实验场景设计以及详细的实验结果分析。

第5章讨论Swift算法的适用边界和未来研究方向。

第6章总结全文并展望未来工作。
